% GENERATED FILE. Edit canonical lessons, then run npm run generate:books.
% canonical-source-hash: fnv1a64:5d286e55679d4cc6
% canonical-lessons: JA-C497-shokuryou, JA-C497-shokuhi, JA-C497-jisui-suru, JA-C497-marumeru, JA-C497-okazu, JA-R497-review-492-497

\chapter{Feeding the Family}
\label{ch:ja-feeding-the-family}
\hlchaptermodality{497}

\begin{chapteropening}
\textbf{By the end of this chapter:} \emph{I can talk about food at home and what it costs.}

Complete the last lesson of chapter 497: I can talk about food at home and what it costs.
\end{chapteropening}

\section[\texorpdfstring{\ja{しょくりょう} (shokuryō)}{しょくりょう (shokuryō)}]{\ja{しょくりょう} (shokuryō) — food supplies; provisions}

\label{lesson:JA-C497-shokuryou}



\begin{quote}
Before the new one: say the Japanese for salt content, then the Japanese for an ingredient.
\end{quote}

\subsection*{You'll want to know: \ja{しょくりょう}}
\textbf{\ja{しょくりょう}} — \emph{shokuryō} — \textquotedblleft{}food supplies; provisions\textquotedblright{}.

\begin{quote}
\textbf{\ja{しょくりょうを} \ja{かばんに} \ja{いれました。}} — \emph{shokuryō o kaban ni iremashita} — \textquotedblleft{}I put the food supplies in the bag.\textquotedblright{}
\end{quote}

\begin{cousinweb}
\textbf{\ja{しょくりょう}} (\emph{shokuryō}) is Sino-Japanese (kango): two Chinese characters read with their on'yomi, readings borrowed from Chinese. They are \emph{shoku}, \textquotedblleft{}eat\textquotedblright{} (Middle Chinese \emph{zyik}) and \emph{ryō}, \textquotedblleft{}materials\textquotedblright{} (\emph{lew}) (KANJIDIC2; Middle Chinese in Baxter's transcription).
\end{cousinweb}

\begin{grammarlens}[title={What you've built}]
One more word for describing what you lived and felt.
\end{grammarlens}

\subsection*{Your turn}
\begin{itemize}
\raggedright
  \item \emph{Say it:} \emph{shokuryō}
  \item \emph{Say it:} \emph{shokuryō}, once more
  \item \emph{Say it:} \emph{shokuryō o kaban ni iremashita} — I put the food supplies in the bag.
  \item \emph{Recall:} say the Japanese for salt content, then the Japanese for an ingredient, then say \emph{shokuryō} again
\end{itemize}

\begin{tcolorbox}[breakable,colback=teal!4,colframe=teal!35!black,title={Before you move on}]
What does \ja{しょくりょう} mean? (\textquotedblleft{}Food supplies; provisions\textquotedblright{}.) Say it once more.
\end{tcolorbox}
\section[\texorpdfstring{\ja{しょくひ} (shokuhi)}{しょくひ (shokuhi)}]{\ja{しょくひ} (shokuhi) — food costs}

\label{lesson:JA-C497-shokuhi}



\begin{quote}
Before the new one: say the Japanese for an ingredient, then the Japanese for food supplies.
\end{quote}

\subsection*{You'll want to know: \ja{しょくひ}}
\textbf{\ja{しょくひ}} — \emph{shokuhi} — \textquotedblleft{}food costs\textquotedblright{}.

\begin{quote}
\textbf{\ja{せんげつは} \ja{しょくひが} \ja{たかかったです。}} — \emph{sengetsu wa shokuhi ga takakatta desu} — \textquotedblleft{}Food costs were high last month.\textquotedblright{}
\end{quote}

\begin{cousinweb}
\textbf{\ja{しょくひ}} (\emph{shokuhi}) is Sino-Japanese (kango): two Chinese characters read with their on'yomi, readings borrowed from Chinese. They are \emph{shoku}, \textquotedblleft{}eat\textquotedblright{} (Middle Chinese \emph{zyik}) and \emph{hi}, \textquotedblleft{}expense\textquotedblright{} (\emph{phj+jH}) (KANJIDIC2; Middle Chinese in Baxter's transcription).
\end{cousinweb}

\begin{grammarlens}[title={What you've built}]
One more word for describing what you lived and felt.
\end{grammarlens}

\subsection*{Your turn}
\begin{itemize}
\raggedright
  \item \emph{Say it:} \emph{shokuhi}
  \item \emph{Say it:} \emph{shokuhi}, once more
  \item \emph{Say it:} \emph{sengetsu wa shokuhi ga takakatta desu} — Food costs were high last month.
  \item \emph{Recall:} say the Japanese for an ingredient, then the Japanese for food supplies, then say \emph{shokuhi} again
\end{itemize}

\begin{tcolorbox}[breakable,colback=teal!4,colframe=teal!35!black,title={Before you move on}]
What does \ja{しょくひ} mean? (\textquotedblleft{}Food costs\textquotedblright{}.) Say it once more.
\end{tcolorbox}
\section[\texorpdfstring{\ja{じすいする} (jisui suru)}{じすいする (jisui suru)}]{\ja{じすいする} (jisui suru) — to cook for oneself}

\label{lesson:JA-C497-jisui-suru}



\begin{quote}
Before the new one: say the Japanese for food supplies, then the Japanese for food costs.
\end{quote}

\subsection*{You'll want to know: \ja{じすいする}}
\textbf{\ja{じすいする}} — \emph{jisui suru} — \textquotedblleft{}to cook for oneself\textquotedblright{}.

\begin{quote}
\textbf{\ja{あには} \ja{まいにち} \ja{じすいしています。}} — \emph{ani wa mainichi jisuishiteimasu} — \textquotedblleft{}My elder brother cooks for himself every day.\textquotedblright{}
\end{quote}

Polite \textbf{\ja{じすいします}}, \emph{jisuishimasu}; a \ja{する} verb. Past (polite): \textbf{\ja{あにが} \ja{じすいしました。}}, \emph{ani ga jisuishimashita}.

\begin{cousinweb}
\textbf{\ja{じすいする}} (\emph{jisui suru}) is the Sino-Japanese (kango) noun \emph{jisui} with the native verb \emph{suru}, \textquotedblleft{}to do\textquotedblright{}: two Chinese characters read with their on'yomi, readings borrowed from Chinese. They are \emph{ji}, \textquotedblleft{}oneself\textquotedblright{} (Middle Chinese \emph{dzijH}) and \emph{sui}, \textquotedblleft{}cook\textquotedblright{} (\emph{tsyhwe}) (KANJIDIC2; Middle Chinese in Baxter's transcription).
\end{cousinweb}

\begin{grammarlens}[title={What you've built}]
One more word for describing what you lived and felt.
\end{grammarlens}

\subsection*{Your turn}
\begin{itemize}
\raggedright
  \item \emph{Say it:} \emph{jisui suru}
  \item \emph{Say it:} \emph{jisui suru}, once more
  \item \emph{Say it:} \emph{ani wa mainichi jisuishiteimasu} — My elder brother cooks for himself every day.
  \item \emph{Recall:} say the Japanese for food supplies, then the Japanese for food costs, then say \emph{jisui suru} again
\end{itemize}

\begin{tcolorbox}[breakable,colback=teal!4,colframe=teal!35!black,title={Before you move on}]
What does \ja{じすいする} mean? (\textquotedblleft{}To cook for oneself\textquotedblright{}.) Say it once more.
\end{tcolorbox}
\section[\texorpdfstring{\ja{まるめる} (marumeru)}{まるめる (marumeru)}]{\ja{まるめる} (marumeru) — to roll into a ball}

\label{lesson:JA-C497-marumeru}



\begin{quote}
Before the new one: say the Japanese for food costs, then the Japanese for to cook for oneself.
\end{quote}

\subsection*{You'll want to know: \ja{まるめる}}
\textbf{\ja{まるめる}} — \emph{marumeru} — \textquotedblleft{}to roll into a ball\textquotedblright{}.

\begin{quote}
\textbf{\ja{ごはんを} \ja{まるめて、おにぎりを} \ja{つくります。}} — \emph{gohan o marumete, onigiri o tsukurimasu} — \textquotedblleft{}I roll the rice into balls to make rice balls.\textquotedblright{}
\end{quote}

Polite \textbf{\ja{まるめます}}, \emph{marumemasu}; an ichidan verb. Past (polite): \textbf{\ja{こどもが} \ja{ねんどを} \ja{まるめました。}}, \emph{kodomo ga nendo o marumemashita}.

\begin{cousinweb}
\textbf{\ja{まるめる}} (\emph{marumeru}) is a native Japanese word (wago), inherited rather than borrowed from Chinese; its kanji was attached later as a kun'yomi, and its older origin is not settled (JMdict, KANJIDIC2).
\end{cousinweb}

\begin{grammarlens}[title={What you've built}]
One more word for describing what you lived and felt.
\end{grammarlens}

\subsection*{Your turn}
\begin{itemize}
\raggedright
  \item \emph{Say it:} \emph{marumeru}
  \item \emph{Say it:} \emph{marumeru}, once more
  \item \emph{Say it:} \emph{gohan o marumete, onigiri o tsukurimasu} — I roll the rice into balls to make rice balls.
  \item \emph{Recall:} say the Japanese for food costs, then the Japanese for to cook for oneself, then say \emph{marumeru} again
\end{itemize}

\begin{tcolorbox}[breakable,colback=teal!4,colframe=teal!35!black,title={Before you move on}]
What does \ja{まるめる} mean? (\textquotedblleft{}To roll into a ball\textquotedblright{}.) Say it once more.
\end{tcolorbox}
\section[\texorpdfstring{\ja{おかず} (okazu)}{おかず (okazu)}]{\ja{おかず} (okazu) — a side dish (eaten with rice)}

\label{lesson:JA-C497-okazu}



\begin{quote}
Before the new one: say the Japanese for to cook for oneself, then the Japanese for to roll into a ball.
\end{quote}

\subsection*{You'll want to know: \ja{おかず}}
\textbf{\ja{おかず}} — \emph{okazu} — \textquotedblleft{}a side dish (eaten with rice)\textquotedblright{}.

\begin{quote}
\textbf{\ja{ばんごはんの} \ja{おかずは} \ja{さかなです。}} — \emph{bangohan no okazu wa sakana desu} — \textquotedblleft{}The side dish for dinner is fish.\textquotedblright{}
\end{quote}

\begin{cousinweb}
\textbf{\ja{おかず}} (\emph{okazu}) is a native Japanese word (wago), inherited rather than borrowed from Chinese, and usually written in kana; its older origin is not settled (JMdict).
\end{cousinweb}

\begin{grammarlens}[title={What you've built}]
One more word for describing what you lived and felt.
\end{grammarlens}

\subsection*{Your turn}
\begin{itemize}
\raggedright
  \item \emph{Say it:} \emph{okazu}
  \item \emph{Say it:} \emph{okazu}, once more
  \item \emph{Say it:} \emph{bangohan no okazu wa sakana desu} — The side dish for dinner is fish.
  \item \emph{Recall:} say the Japanese for to cook for oneself, then the Japanese for to roll into a ball, then say \emph{okazu} again
\end{itemize}

\begin{tcolorbox}[breakable,colback=teal!4,colframe=teal!35!black,title={Before you move on}]
What does \ja{おかず} mean? (\textquotedblleft{}A side dish (eaten with rice)\textquotedblright{}.) Say it once more.
\end{tcolorbox}
\section[(dialogue)]{Looking back: from Simmer, Soak and Steam to Feeding the Family}

\label{lesson:JA-R497-review-492-497}



\begin{quote}
Say the Japanese for to be cooked through and to boil.
\end{quote}

\subsection*{Your turn}
Say these aloud:

\begin{itemize}
\raggedright
  \item to be cooked through, to boil, to dip in liquid, to steam (food), to peel
  \item a freezer, a cupboard, a rice cooker, the washing-up, to freeze (food)
  \item an onion, a cucumber, a mushroom, a bamboo shoot, a shiitake mushroom
  \item a smell, a fragrance, flavour, fresh, light (of food)
  \item soft, salty, raw, salt content, an ingredient
  \item food supplies, food costs, to cook for oneself, to roll into a ball, a side dish (eaten with rice)
\end{itemize}

\begin{tcolorbox}[breakable,colback=teal!4,colframe=teal!35!black,title={Before you move on}]
To be cooked through? (\textbf{\ja{にえる}}, \emph{nieru}.) To boil? (\textbf{\ja{わく}}, \emph{waku}.) To dip in liquid? (\textbf{\ja{ひたす}}, \emph{hitasu}.) To steam (food)? (\textbf{\ja{ふかす}}, \emph{fukasu}.) To peel? (\textbf{\ja{むく}}, \emph{muku}.) A freezer? (\textbf{\ja{れいとうこ}}, \emph{reitōko}.) A cupboard? (\textbf{\ja{とだな}}, \emph{todana}.) A rice cooker? (\textbf{\ja{すいはんき}}, \emph{suihanki}.) The washing-up? (\textbf{\ja{あらいもの}}, \emph{araimono}.) To freeze (food)? (\textbf{\ja{れいとうする}}, \emph{reitō suru}.) An onion? (\textbf{\ja{たまねぎ}}, \emph{tamanegi}.) A cucumber? (\textbf{\ja{きゅうり}}, \emph{kyūri}.) A mushroom? (\textbf{\ja{きのこ}}, \emph{kinoko}.) A bamboo shoot? (\textbf{\ja{たけのこ}}, \emph{takenoko}.) A shiitake mushroom? (\textbf{\ja{しいたけ}}, \emph{shiitake}.) A smell? (\textbf{\ja{におい}}, \emph{nioi}.) A fragrance? (\textbf{\ja{かおり}}, \emph{kaori}.) Flavour? (\textbf{\ja{ふうみ}}, \emph{fūmi}.) Fresh? (\textbf{\ja{しんせん}}, \emph{shinsen}.) Light (of food)? (\textbf{\ja{あっさり}}, \emph{assari}.) Soft? (\textbf{\ja{やわらかい}}, \emph{yawarakai}.) Salty? (\textbf{\ja{しおからい}}, \emph{shiokarai}.) Raw? (\textbf{\ja{なま}}, \emph{nama}.) Salt content? (\textbf{\ja{えんぶん}}, \emph{enbun}.) An ingredient? (\textbf{\ja{しょくざい}}, \emph{shokuzai}.) Food supplies? (\textbf{\ja{しょくりょう}}, \emph{shokuryō}.) Food costs? (\textbf{\ja{しょくひ}}, \emph{shokuhi}.) To cook for oneself? (\textbf{\ja{じすいする}}, \emph{jisui suru}.) To roll into a ball? (\textbf{\ja{まるめる}}, \emph{marumeru}.) A side dish (eaten with rice)? (\textbf{\ja{おかず}}, \emph{okazu}.)
\end{tcolorbox}
